% ===================================================================
%  TABELUL Nr. 12 — Gara. — Căile ferate. — Vasele.
%  Traduction 2026 d'apres texto/io, controlee sur texto/fr, texto/en
%  et texto/es. Consigne : voir texto/ro/10-tabelul-01.tex.
%
%  MEME OBSERVATION QU'AU TABLEAU 10 : « vagon », « peron »,
%  « semafor », « tender », « macaz » — le rail est arrive tout fait.
%  Ce qui reste roumain est ce que le rail a trouve en place :
%  « hamal » le porteur, qui est turc, et « geamantan » la valise, qui
%  l'est aussi.
% ===================================================================

%%K t12-apar-1 apar
\VUsaut{4.0mm}
\VUtitre{15.0pt}{60}{TABELUL Nr. 12}
\VUsaut{2.4mm}
\VUpk{11.4pt}{40}{Gara. --- Căile ferate. --- Vasele.}
\VUsaut{3.4mm}
\textsc{Observație.} --- \textit{Orarele sunt în fiecare țară altele, de aceea am
luat ca pildă orele} \VUgras{tabelului francez}.

%%K t12-01-1 p
1. --- Am făcut de curând o călătorie prin Germania. Înainte de plecare am cumpărat
un \VUgras{mersul trenurilor} \textsuperscript{(15)}, ca să aflu ceasurile de
plecare. Încredințându-mă că trenul cel mai potrivit pleacă din Paris la nouă seara
și ajunge la Strasbourg la unu și treisprezece, mi-am pregătit călătoria, iar a
doua zi m-am lăsat dus la gară.

%%K t12-02-1 p
2. --- Când am intrat în sala mare de plecare, în urma unui bătrân \VUgras{preot}
de țară \textsuperscript{(2)}, care își ducea \VUgras{sacul de călătorie}
\textsuperscript{(3)} în mână, se strânseseră deja mulți \VUgras{călători}
\textsuperscript{(1)}.

%%K t12-02-2 p
Fiindcă aveam de trimis o \VUgras{telegramă} \textsuperscript{(5)}, l-am rugat pe un
\VUgras{controlor} \textsuperscript{(6)} să îmi arate \VUgras{telegraful}
\textsuperscript{(4)}.

%%K t12-03-1 p
3. --- Înainte de a intra în \VUgras{sala de așteptare} \textsuperscript{(7)}, am
mers să îmi iau un \VUgras{bilet} dus și întors \textsuperscript{(9)}.

%%K t12-04-1 p
4. --- La \VUgras{ghișeu} \textsuperscript{(8)} nu am așteptat mult, fiindcă erau
puțini oameni. Un \VUgras{soldat} \textsuperscript{(10)}, care mergea în permisie,
a fost atât de binevoitor încât mi-a lăsat rândul lui, așa că am avut destul timp
să cer câteva lămuriri de la \VUgras{șeful gării} \textsuperscript{(13)}, care
stătea de vorbă cu \VUgras{comisarul} \textsuperscript{(12)} de pază și cu
\VUgras{jandarmul} de serviciu \textsuperscript{(11)}.

%%K t12-tit-1 sub
\VUpk{10.2pt}{40}{Predarea bagajelor.}
\VUsaut{3.0mm}

%%K t12-05-1 p
5. --- După ce am cumpărat un ziar la \VUgras{librărie} \textsuperscript{(14)}, am
pus să mi se cântărească \VUgras{bagajul} \textsuperscript{(17)}, pe care un
\VUgras{hamal} \textsuperscript{(16)} tocmai îl adusese pe un \VUgras{cărucior de
bagaje} \textsuperscript{(19)}. \VUgras{Funcționarul} \textsuperscript{(22)} de la
\VUgras{cântar} \textsuperscript{(21)} mi-a spus că lada mea cântărește treizeci și
cinci de kilograme; asta făcea un \VUgras{surplus} de cinci kilograme, pentru care
am plătit un adaos la \VUgras{ghișeul} bagajelor \textsuperscript{(23)}.

%%K t12-06-1 p
6. --- Fiindcă venisem prea devreme, am avut vreme să privesc mișcarea călătorilor
prin gară. Unii își predau sau își ridicau bagajele de mână la \VUgras{casa de
bagaje} \textsuperscript{(25)}; alții cercetau \VUgras{mersul trenurilor}
\textsuperscript{(26)}. Călătorii care soseau se grăbeau spre \VUgras{ieșire}
\textsuperscript{(32)} cu \VUgras{geamantanul} \textsuperscript{(24)} în mână.
Unii așteptau la \VUgras{eliberarea bagajelor} \textsuperscript{(27)} și își arătau
\VUgras{recipisa} \textsuperscript{(29)}, ca să își primească lăzile,
\VUgras{cutiile} \textsuperscript{(28)} sau \VUgras{coșurile}
\textsuperscript{(30)}.

%%K t12-07-1 p
7. --- \VUgras{Funcționarii de acciz} \textsuperscript{(33)} cercetau cu grijă
pachetele, ca să afle dacă este ceva de declarat.

%%K t12-08-1 p
8. --- Unii călători se îndreptau spre \VUgras{bufet} \textsuperscript{(34)}, unde
un \VUgras{chelner} \textsuperscript{(35)} se silea să îi servească. Trebuiau să se
grăbească, fiindcă \VUgras{trenul} \textsuperscript{(36)}, care este accelerat, nu
avea să stea mult în gară.

%%K t12-09-1 p
9. --- Apoi am ieșit pe peronul de plecare și am căutat un \VUgras{vagon} direct
\textsuperscript{(37)}, ca să nu fiu silit să schimb pe drum. Am intrat într-un
vagon cu coridor, care are un \VUgras{compartiment} \textsuperscript{(38)} pentru
fumători și unul păstrat pentru «doamne care călătoresc singure».

%%K t12-tit-2 sub
\VUpk{10.2pt}{40}{Plecare noaptea.}
\VUsaut{3.0mm}

%%K t12-10-1 p
10. --- După ce am urcat pe \VUgras{scară} \textsuperscript{(40)}, am închis
\VUgras{ușița} \textsuperscript{(39)}, am ridicat \VUgras{storul}
\textsuperscript{(41)} și mi-am pus bagajele de mână în \VUgras{plasă}
\textsuperscript{(43)}, m-am așezat pe \VUgras{banchetă} \textsuperscript{(42)}.
Când l-am văzut pe \VUgras{lampagiu} \textsuperscript{(44)} aprinzând lămpile,
mi-am adus aminte că voi petrece o noapte în vagon și l-am chemat pe funcționarul
care se îngrijește de închirierea \VUgras{pernelor} \textsuperscript{(45)}, ca să
cer una.

%%K t12-11-1 p
11. --- Conductorul trenului a venit să verifice biletele. L-am întrebat de ce nu am
plecat încă, fiindcă este ora potrivită. Mi-a răspuns că \VUgras{șeful de tren}
\textsuperscript{(46)} pune să fie încărcate biciclete în \VUgras{vagonul de
bagaje} \textsuperscript{(47)}. Pe deasupra nu fuseseră schimbate încă toate
\VUgras{încălzitoarele de picioare} \textsuperscript{(48)}.

%%K t12-12-1 p
12. --- Dar aveam să plecăm curând, fiindcă \VUgras{fochistul}
\textsuperscript{(50)} lua cărbune pe \VUgras{tenderul} său \textsuperscript{(49)},
ca să ațâțe focul \VUgras{locomotivei} \textsuperscript{(54)}, iar
\VUgras{mecanicul} \textsuperscript{(51)} aștepta numai semnul \VUgras{ajutorului de
șef de gară} \textsuperscript{(52)}, care stătea pe \VUgras{peron}
\textsuperscript{(53)}.

%%K t12-13-1 p
13. --- \VUgras{Cazanul} \textsuperscript{(56)} locomotivei duduia surd;
\VUgras{coșul} \textsuperscript{(55)} arunca șuvoaie de fum amestecat cu
\VUgras{abur} \textsuperscript{(60)}. S-a auzit un \VUgras{fluierat} ascuțit
\textsuperscript{(59)} și \VUgras{pistoanele} \textsuperscript{(58)} au început să
se miște, punând în mers \VUgras{roțile} \textsuperscript{(57)} mașinii grele.

%%K t12-tit-3 sub
\VUsaut{3.0mm}
\VUpk{10.2pt}{40}{Am pornit!}
\VUsaut{3.0mm}

%%K t12-14-1 p
14. --- Trenul, mergând pe \VUgras{șine} \textsuperscript{(64)}, își sporea viteza;
\VUgras{peroanele} \textsuperscript{(65)} s-au sfârșit; am trecut de
\VUgras{closete} \textsuperscript{(66)}, și de asemenea de un șir de vagoane oprite
pe \VUgras{linia} cealaltă \textsuperscript{(63)}: \VUgras{vagoane de dormit}
\textsuperscript{(67)}, un \VUgras{vagon-restaurant} \textsuperscript{(68)}, un
\VUgras{vagon poștal} \textsuperscript{(69)}.

%%K t12-noto-1 noto
\VUnotes{143.2mm}{%
\textit{(*) Observație.} --- \textit{Călătoria este, firește, descrisă așa cum se
petrecea la hotarul de dinainte de război.}%
}

%%K t12-15-1 p
15. --- Apoi ne-a trecut prin fața ochilor \VUgras{gara de mărfuri}
\textsuperscript{(71)}. Sub șoproane funcționarii încărcau \VUgras{mărfurile}
\textsuperscript{(72)} trimise cu viteză mică. \VUgras{Manevranții}
\textsuperscript{(76)} de lângă \VUgras{turnul de apă} \textsuperscript{(73)}
mutau \VUgras{vagoane de vite} \textsuperscript{(75)} și \VUgras{platforme}
\textsuperscript{(79)}, ca să alcătuiască un \VUgras{tren de marfă}
\textsuperscript{(80)}.

%%K t12-16-1 p
16. --- Am trecut de ultimul \VUgras{macaz} \textsuperscript{(78)}, lângă care
stătea acarul; am trecut de \VUgras{semafor} \textsuperscript{(86)}, și gara a
pierit în depărtare.

%%K t12-17-1 p
17. --- Curând am ieșit pe un \VUgras{pod de fier} \textsuperscript{(74)}, care
trece peste \VUgras{râu} \textsuperscript{(81)}. După ce am trecut acest pod, am
mers de-a lungul râului, pe care \VUgras{remorchere} puternice
\textsuperscript{(82)} trăgeau \VUgras{șlepuri} grele \textsuperscript{(83)}. Am
văzut și un \VUgras{vapor de călători} \textsuperscript{(84)}, care tocmai se
apropia de un \VUgras{ponton} \textsuperscript{(85)}.

%%K t12-18-1 p
18. --- După ce am trecut cu iuțeală nebună prin câteva gări, am străbătut câteva
\VUgras{tuneluri} \textsuperscript{(87)} și am lăsat în urmă nenumăratele
\VUgras{căsuțe} \textsuperscript{(88)} ale \VUgras{cantonierilor}. Legănat de
zdruncinăturile trenului, am adormit adânc.

%%K t12-tit-4 sub
\VUsaut{3.0mm}
\VUpk{10.2pt}{40}{Gara Deutsch-Avricourt.}
\VUsaut{3.0mm}

%%K t12-19-1 p
19. --- M-a trezit deodată glasul unui funcționar care striga: «Deutsch-Avricourt!
\textsuperscript{(*)} Toată lumea coboară!» A fost cercetarea vămii și toate
formalitățile plicticoase ale hotarului. A trebuit să îmi potrivesc ceasul după
\VUgras{ceasul cel mare} al gării \textsuperscript{(89)}, care mergea cu cincizeci
și cinci de minute înainte; abia treaz, deosebeam cu greu \VUgras{limba mare}
\textsuperscript{(90)} de \VUgras{cea mică} \textsuperscript{(91)}; însuși
\VUgras{cadranul} \textsuperscript{(92)} era cu totul tulburat de ceața dimineții.

%%K t12-20-1 p
20. --- În vreme ce vameșii își făceau cercetarea, căscând, deslușeam
\VUgras{afișele} \textsuperscript{(93)} care împodobesc zidurile gării. Am luat
gustarea ca să treacă vremea, am cercetat harta \VUgras{rețelei} germane
\textsuperscript{(94)}, ca să văd cât mai am până la Strasbourg, și m-am întors în
vagonul meu, întărit de nădejdea că voi ajunge curând la ținta călătoriei.
\VUsaut{6.0mm}
\VUfilet{22mm}
